\scanpage{30}
Revenant au cas $W$ une sphère, on trouve
\[\left\{\begin{array}{ll}H^i(W-W_0)=0&\text{si }i>n+1\\H^i(W_0)=0&\text{si }i\ne n-1,n,0,2n-1\end{array}\right.\]
\margin{[en fait, si $i\ne\deleted{\unsure{n,n+1}}\deleted{\unsure{]}}0,1,n,n+1$]}
\margin{(d'abord si $i\ge n+1$ et $i\ne2n-1$, puis le reste par dualité]}

La suite exacte associée à $d)$ est triviale, sauf la portion
\[\left\{\begin{gathered}
0\longrightarrow H^0(W)\longrightarrow H^0(W-W_0)\longrightarrow0\\
0\longrightarrow H^n(W-W_0)\longrightarrow H^{n-1}(W_0)\longrightarrow0\\
0\longrightarrow H^{n+1}(W-W_0)\longrightarrow H^n(W_0)\xrightarrow[\text{si }n\ne1]{}0\\
0\longrightarrow H^{2n-1}(W_0)\longrightarrow H^{2n+1}(W)\longrightarrow0
\end{gathered}\right.\]
\[0\longrightarrow H^1(W-W_0)\longrightarrow H^0(W_0)\longrightarrow0\]
\deleted{\[\unsure{n+1\ne2n+1}\]}
\deleted{\[n+2\ne2n+1\]}

[Si $n=1$, on trouve
\[\begin{gathered}
0\longrightarrow H^0(W)\xrightarrow{\sim}H^0(W-W_0)\mathbin{-}0\\
0\longrightarrow H^1(W-W_0)\longrightarrow H^0(W_0\deleted{\unsure{\Lambda}})\mathbin{-}0\\
0\longrightarrow H^2(W-W_0)\longrightarrow H^1(W_0)\longrightarrow\underset{\Vert\;\Lambda}{H^3(W)}\longrightarrow0
\end{gathered}\]
]
\figure{\begin{tikzpicture}\draw (0,0)--(13,0.6);\end{tikzpicture}}
Notons que, par \unsure{passage} à la limite, \deleted{\unsure{les les}} les relations
\[H^i(W\unsure{-}W_0,\Lambda)=0\quad\text{pour }i>n\deleted{\unsure{+1}}\]
donnent
\[\boxed{H^i(M,\Lambda)=0\quad\text{si }i>n\deleted{\unsure{+1}}}\]
\margin{[\unsure{Indépendt} du fait que $W$ non singulier] \inserted{et \unsure{même} que $W_0$ non singulier}}

La suite exacte relative (si $W$ non singulier)
\[\mathbin{-}\longrightarrow H^{\unsure{2n-i}}(M)^\vee\longrightarrow H^i(M)\longrightarrow H^i(W_0)\longrightarrow H^{2n-i\unsure{+1|-1}}(M)^\vee\longrightarrow H^{i+1}(M)\longrightarrow H^{\unsure{i+1}}(W_0)\]
montre donc que

\deleted{Si $i>n+1$}
\mspage{29}
